\documentclass[t,xcolor=pst,dvips]{beamer}
\mode<presentation>

\usepackage{etex}

\usetheme{Madrid}
% other themes: Warsaw, AnnArbor, Antibes, Bergen, Berkeley, Berlin, Boadilla, boxes, CambridgeUS, Copenhagen, Darmstadt, default, Dresden, Frankfurt, Goettingen,
% Hannover, Ilmenau, JuanLesPins, Luebeck, Madrid, Maloe, Marburg, Montpellier, PaloAlto, Pittsburg, Rochester, Singapore, Szeged, classic

\setbeamertemplate{navigation symbols}{\insertslidenavigationsymbol}

\usecolortheme{dolphin}
%\usecolortheme{seagull}
% color themes: albatross, beaver, beetle, crane, default, dolphin, dov, fly, lily, orchid, rose, seagull, seahorse, sidebartab, structure, whale, wolverine

%\usefonttheme{serif}
% font themes: default, professionalfonts, serif, structurebold, structureitalicserif, structuresmallcapsserif

% pdf is displayed in full screen mode automatically
%\hypersetup{pdfpagemode=FullScreen}

%\AtBeginSection[]
%{
%  \begin{frame}<beamer>
%    \frametitle{Outline}
%    \tableofcontents[currentsection,currentsubsection]
%  \end{frame}
%}

% define your own colours:
\definecolor{Red}{rgb}{1,0,0}
\definecolor{Blue}{rgb}{0,0,1}
\definecolor{Green}{rgb}{0,1,0}
\definecolor{magenta}{rgb}{1,0,.6}
\definecolor{lightblue}{rgb}{0,.5,1}
\definecolor{lightpurple}{rgb}{.6,.4,1}
\definecolor{gold}{rgb}{.6,.5,0}
\definecolor{orange}{rgb}{1,0.4,0}
\definecolor{hotpink}{rgb}{1,0,0.5}
\definecolor{newcolor2}{rgb}{.5,.3,.5}
\definecolor{newcolor}{rgb}{0,.3,1}
\definecolor{newcolor3}{rgb}{1,0,.35}
\definecolor{darkgreen1}{rgb}{0, .35, 0}
\definecolor{darkgreen}{rgb}{0, .6, 0}
\definecolor{darkred}{rgb}{.75,0,0}

\xdefinecolor{olive}{cmyk}{0.64,0,0.95,0.4}
\xdefinecolor{purpleish}{cmyk}{0.75,0.75,0,0}

%\usepackage{beamerinnerthemerounded}
% inner themes include circles, default, inmargin, rectangles, rounded

%\usepackage{beamerouterthemesmoothbars}
% outer themes include default, infolines, miniframes, shadow, sidebar, smoothbars, smoothtree, split, tree

\useoutertheme[subsection=false]{smoothbars}

% to have the same footer on all slides
\setbeamertemplate{footline}[text line]{
\includegraphics[height=15pt]{logo-high-res.eps}\hfill 
\raisebox{5pt}{Math 314:  Statistics}\hfill 
\raisebox{5pt}{Chapter 2:  Observational Studies}\hfill
\raisebox{5pt}{\insertframenumber/\pageref{lastpage}}}
%\setbeamertemplate{footline}[text line]{} % or empty footer

% include packages
\usepackage{subfigure}
\usepackage{multicol}
\usepackage{amsmath}
\usepackage{epsfig}
\usepackage{graphicx}
\usepackage[all,knot]{xy}
\xyoption{arc}
\usepackage{url}
\usepackage{multimedia}
\usepackage{hyperref}
%\usepackage{setspace}
\usepackage{pstricks,pst-grad,pst-text,pst-node,multido,pst-plot,calc,pst-3dplot}

\title{Math 314:  Statistics}
\subtitle{Chapter 2:  Observational Studies}
\author{Ralph Wojtowicz}
\institute{CME Department\\ Shepherd University}
%\date{\scriptsize 24 August 2011}


\begin{document}

%\frame[plain]{
%	\titlepage
%}

%\section[Outline]{}


\begin{frame}[plain]
\definecolor{myblue}{rgb}{0,0,0.6}
\definecolor{grayA}{rgb}{0.95,0.95,0.95}
\definecolor{grayB}{rgb}{0.98,0.98,0.98}
\begin{center}

\newcommand{\IMG}{
\scalebox{0.63}{\begin{pspicture}(-2,-2)(2,2)
\psframe(-1.5,-0.5)(8.5,5.25)
\psset{xunit=1.4,yunit=1.2,nodesep=2pt}
\rput(4,4){\rnode{Studies}{Studies}}
\rput(5,3){\rnode{NoControls}{No Controls}}
\rput(3,3){\rnode{Controls}{Controls}}
\rput(4,2){\rnode{Historical}{Historical}}
\rput(2,2){\rnode{Contemporaneous}{Contemporaneous}}
\rput(1,1){\rnode{Controlled}{\begin{tabular}{c}Controlled\\ Experiment\end{tabular}}}
\rput(3,1){\rnode{Observational}{\begin{tabular}{c}Observational\\ Studies\end{tabular}}}
\rput(0,0){\rnode{Randomized}{Randomized}}
\rput(2,0){\rnode{NotRandomized}{Not Randomized}}
\ncline{Studies}{Controls}\ncline{Studies}{NoControls}
\ncline{Controls}{Contemporaneous}\ncline{Controls}{Historical}
\ncline{Contemporaneous}{Controlled}\ncline{Contemporaneous}{Observational}
\ncline{Controlled}{Randomized}\ncline{Controlled}{NotRandomized}
\end{pspicture}}}

%\begin{pspicture}(0,0)(7,4.8)
\begin{pspicture}(-6,-7)(6,2)
\psframe[linewidth=0.02,linecolor=gray](-6.2,-7)(6.2,2.2)
\psframe[linewidth=0.02,linecolor=gray](-6.15,-6.95)(6.15,2.15)
\rput(0,1.2){\color{myblue}\large Math 314:  Statistics}
\rput(0,0.2){\color{myblue}Chapter 2:  Observational Studies}
\rput(-2.3,-3.5){\IMG}
%\psframebox(0,0)(4,4)
\rput(0,-4.2){\scriptsize Dr.~Ralph Wojtowicz}
\rput(0,-4.7){\scriptsize CME Department}
\rput(0,-5.4){\includegraphics[height=1cm]{logo-high-res.eps}}
%
%\rput(0,-6.5){\scriptsize 24 August 2011}
\end{pspicture}
\end{center}

\end{frame}










\addtocounter{page}{-1}
\addtocounter{framenumber}{-1}
\frame{\tableofcontents}

\section{Introduction}
\subsection{Introduction}
\begin{frame}[t]\frametitle{Introduction}
{\small 
\begin{itemize}
\item \textbf{Controlled experiment}: Investigators decide who will be in the treatment group and who will 
  be in the control group
     \begin{itemize}
     \item Example:  Salk vaccine trials discussed in Chapter 1
     \item Example:  Coronary Drug Project discussed in Section 2.2
     \item The control and treatment populations are similar --- except in the application of treatment
     \end{itemize}
\item \textbf{Observational study}: Subjects assign themselves to the groups
     \begin{itemize}
     \item Example:  Smoking studies
%     \item The two populations differ in ways other than the treatment
     \item \textbf{Association} between treatment and outcome is circumstantial evidence for causation.    
     \item Association does not prove causation.  \textbf{Confounding} factors may exist.
     \item Observational studies can be powerful tools but can also be misleading.% due to confounding.  
       \begin{itemize}
         \item Were the control and treatment groups similar?
         \item Did the two populations differ in ways other than the treatment?
       \end{itemize}
     \item Technique:  compare small, more homogeneous groups (e.g.,~age, sex)
     \end{itemize}
\end{itemize}}
\end{frame}

\subsection[Studies]{Classification of Studies}
\begin{frame}[t]\frametitle{Classification of Studies}
{\footnotesize
\begin{center}
\begin{pspicture}(0,0)(7,4.8)
\psset{xunit=1.4,yunit=1.2,nodesep=2pt}
\rput(4,4){\rnode{Studies}{Studies}}
\rput(5,3){\rnode{NoControls}{No Controls}}
\rput(3,3){\rnode{Controls}{Controls}}
\rput(4,2){\rnode{Historical}{Historical}}
\rput(2,2){\rnode{Contemporaneous}{Contemporaneous}}
\rput(1,1){\rnode{Controlled}{\begin{tabular}{c}Controlled\\ Experiment\end{tabular}}}
\rput(3,1){\rnode{Observational}{\begin{tabular}{c}Observational\\ Studies\end{tabular}}}
\rput(0,0){\rnode{Randomized}{Randomized}}
\rput(2,0){\rnode{NotRandomized}{Not Randomized}}
\ncline{Studies}{Controls}\ncline{Studies}{NoControls}
\ncline{Controls}{Contemporaneous}\ncline{Controls}{Historical}
\ncline{Contemporaneous}{Controlled}\ncline{Contemporaneous}{Observational}
\ncline{Controlled}{Randomized}\ncline{Controlled}{NotRandomized}
%\psframebox(0,0)(4,4)
\end{pspicture}
\end{center}}
\bigskip

{\small See page 28. Compare the figure on page 7.}
\end{frame}

\section{The Clofibrate Trial}
\subsection{Clofibrate}
\begin{frame}[t]\frametitle{Clofibrate}
{\small
\begin{itemize}
\item Coronary Drug Project:  randomized, controlled double-blind experiment 
  (placebo $=$ lactose) to evaluate heart attack prevention drugs
\item 8,341 patients followed for five years (5,552 got treatment, 2,789 controls)
\item Clofibrate:  a cholesterol reduction drug evaluated in the study
\item Comparing 20\% to 21\% shows that clofibrate did not save lives.
\item Many subjects did not take their medicine (non-adherers).
\item Compare 15\% to 15\% (not to 21\% or 25\%)  to control for adherence.
\end{itemize}}\vspace{-8pt}
%
\begin{center}
{\newcommand{\Z}{\hphantom{0,}}\footnotesize\begin{tabular}{lcccc}
& \Z\Z Clofibrate\span  & \Z\Z Placebo\span\\[2pt]\hline
         & Number & Deaths & Number & Deaths\vphantom{\Large Y}\\
Adherers     & \Z708 & 15\% & 1,813 & 15\%\\
Non-adherers & \Z357 & 25\% & \Z882 & 28\%\\
Total group  & 1,103 & 20\% & 2,789 & 21\%
\end{tabular}} \vspace{-8pt}
\end{center}
{\small
\begin{itemize}
\item Conclusions:  (i) Clofibrate does not have an effect.  (ii) Adherers are different from non-adherers.
\end{itemize}}
\end{frame}

\section{More Examples}
\subsection{Confounding Factors and Associations}
\begin{frame}[t]\frametitle{Confounding Factors and Associations}
{\small
\begin{itemize}
\item Pellagra:  Among many associations between the 18${}^{th}$ century disease and other factors, lack of niacin
  was found to be the underlying cause.
\item Cervical Cancer and Circumcision: Human papilloma virus was found to be the underlying cause
  and explained differences in cancer rates between particular populations in the 1950s.
\item Ultrasound and low birthweight:  The confounding factor of problem pregnancy was found to 
  explain an association between ultrasound and low birthweight.  Randomized controlled experiments 
  showed that ultrasounds may be protective.
\item The Samaritans and Suicide:  A confounding factor explained an association between the 
  expansion of a volunteer welfare organization and a decrease in the English suicide rate 
  in 1964--1970.
\end{itemize}
}
\end{frame}


\section{Sex Bias in Graduate Admissions}
\subsection{Sex Bias in Graduate Admissions}
\begin{frame}[t]\frametitle{}
{\small
\begin{itemize}
\item Observational study on sex bias in admissions at UC, Berkeley in 1973
   \begin{itemize}
   \item 44\% of 8,442 male applicants were admitted
   \item 35\% of 4,321 female applicants were admitted
   \end{itemize}
\item Compare admissions to the six largest majors:\vspace{-9pt}
\end{itemize}}

{\footnotesize
\begin{center}
\newcommand{\Z}{\hphantom{0}}
\begin{tabular}{ccccc}
      & \Z\Z\Z\Z\Z\Z Men\span   & \Z\Z\Z\Z\Z Women\span\\[3pt]
      & Number of  & Percent  & Number of  & Percent\\
Major & Applicants & Admitted & Applicants & Admitted\\[2pt]\hline
A     & {\color{red}\textbf{825}}        
                   & {\color{red}\textbf{62}}       
                              & {\color{red}\textbf{108}}        
                                           & {\color{red}\textbf{82}}\vphantom{\large Y}\\
B     & 560        & 63       & 25         & 68\\
C     & 325        & 37       & 593        & 34\\
D     & 417        & 33       & 375        & 35\\
E     & {\color{blue}\textbf{191}}        
                   & {\color{blue}\textbf{28}}       
                              & {\color{blue}\textbf{393}}        
                                           & {\color{blue}\textbf{24}}\\
F     & 373        & \Z6      & 341        & \Z7
\end{tabular}\vspace{-16pt}
\end{center}}
%
{\small
\begin{itemize}
\item Major A:  Less selective but few women and many men applied
\item Major E:  Highly-selective but many women and few men applied
\item Simpson's paradox:  Relationships between percentages in subgroups can 
  be reversed when the subgroups are combined.
\end{itemize}}

\end{frame}

\section{Confounding}
\subsection{Confounding}
\begin{frame}[t]\frametitle{Confounding}

{\small
\begin{itemize}
\item \textbf{Confounding}:  A difference between the treatment and control groups --- other than the treatment ---
   that affects the responses being studied
\item Confounders must be assocated with both:
   \begin{itemize} 
     \item The disease/outcome and 
     \item The exposure/treatment.
   \end{itemize}
\item Hidden confounders are a major problem in observational studies.
\item Examples:
   \begin{itemize}
   \item NFIP polio vaccine study:  family income
   \item Portacaval shunt studies:  health of patients selected for surgery
   \item Coronary bypass surgery studies:  health of patients selected for surgery
   \item Cervical cancer study: sexual activity
   \end{itemize}
\end{itemize}}

\label{lastpage}
\end{frame}

\end{document} 





\begin{pspicture}(-6,-7)(6,2)
\psframe[linewidth=0.02,linecolor=gray](-6.2,-7)(6.2,2.2)
\psframe[linewidth=0.02,linecolor=gray](-6.15,-6.95)(6.15,2.15)
\rput(0,1.2){\color{myblue}\large Math 314:  Statistics}
\rput(0,0.2){\color{myblue}Chapter 2:  Observational Studies}
%
\psframe[linewidth=0.02,fillstyle=solid,fillcolor=grayA](-2,-3)(2,-1)
  \rput[tl](-1.9,-1.1){\tiny Population (parameters)}
\psframe[linewidth=0.02,fillstyle=solid,fillcolor=grayB](-0.25,-2.75)(1.75,-2)
  \rput[tl](-0.15,-2.1){\tiny Sample (statistics)}
%
\rput(0,-4.2){\scriptsize Dr.~Ralph Wojtowicz}
\rput(0,-4.7){\scriptsize CME Department}
\rput(0,-5.4){\includegraphics[height=1cm]{logo-high-res.eps}}
%
\rput(0,-6.5){\scriptsize 24 August 2011}
\end{pspicture}
